% (c) Nikita Lisitsa, lisyarus@gmail.com, 2026

\documentclass[10pt]{beamer}

\usepackage[T2A]{fontenc}
\usepackage[russian]{babel}
\usepackage{minted}

\usepackage[most]{tcolorbox}
\tcbuselibrary{minted}

\usepackage{graphicx}
\graphicspath{ {./images/} }

\usetheme{metropolis}

\usepackage{adjustbox}

\usepackage{color}
\usepackage{soul}

\usepackage{hyperref}

\definecolor{codebg}{RGB}{29,35,49}

\setminted{fontsize=\footnotesize}

\setlength\partopsep{-\topsep}

\definecolor{blue}{rgb}{0,0,1}
\definecolor{red}{rgb}{1,0,0}

\usemintedstyle{tango}

\makeatletter
\newcommand{\slideimage}[1]{
  \begin{figure}
    \begin{adjustbox}{width=\textwidth, totalheight=\textheight-2\baselineskip-2\baselineskip,keepaspectratio}
      \includegraphics{#1}
    \end{adjustbox}
  \end{figure}
}
\makeatother

\title{Компьютерная графика}
\subtitle{Практика 8: Shadow mapping}
\date{2026}

\setbeamertemplate{footline}[frame number]

\begin{document}

\frame{\titlepage}

\begin{frame}[fragile]
\frametitle{Практика 8}
\begin{itemize}
\item Камеру можно двигать и вращать стрелками
\end{itemize}
\end{frame}

\begin{frame}[fragile]
\frametitle{Практика 8}
\slideimage{0.png}
\end{frame}

\begin{frame}[fragile]
\frametitle{Задание 1}
Выводим дебажный прямоугольник для shadow map
\begin{itemize}
\item В практике уже есть \mintinline{cpp}|debug_depth_texture_pipeline|
\item Используем её, чтобы нарисовать дебажный \textbf{квадрат} в левом нижнем углу
\item Для этого нужно правильно выставить \mintinline{cpp}|DebugDepthTextureImmediates|
\item \mintinline{cpp}|draw(vertex_count = 6)|
\item \textbf{NB:} квадрат должен оставаться квадратом при изменении размеров окна!
\end{itemize}
\end{frame}

\begin{frame}[fragile]
\frametitle{Задание 1}
\slideimage{1.png}
\end{frame}

\begin{frame}[fragile]
\frametitle{Задание 2}
Создаём shadow map и sampler
\begin{itemize}
\item Создаём текстуру shadow map: \(1024\times 1024\), \mintinline{cpp}$usage = RENDER_ATTACHMENT | TEXTURE_BINDING$, \mintinline{cpp}|format = Depth24Plus|
\item Создаём sampler: \mintinline{cpp}|address_mode = CLAMP_TO_EDGE|, \mintinline{cpp}|filter = NEAREST|
\item Добавляем их в lights bind group как binding 1 и 2, соответственно
\item Для bind group layout: \mintinline{cpp}|texture.sample_type = Depth|, \mintinline{cpp}|sampler = Filtering|, 
\item \textbf{NB:} Пока ничего не изменится, так как мы не используем shadow map
\end{itemize}
\end{frame}

\begin{frame}[fragile]
\frametitle{Задание 3}
Читаем shadow map в дебажном квадрате
\begin{itemize}
\item Добавляем shadow map и sampler как bindings в \mintinline{text}|debug_depth_texture_shader.wgsl|
\item Для shadow map тип -- \mintinline{wgsl}|texture_depth_2d|
\item Во фрагментном шейдере читаем shadow map \mintinline{wgsl}|textureSample(shadowMap, shadowMapSampler, in.texcoord)| и выводим его как цвет
\item \textbf{NB:} Пока там записаны нулевые/мусорные значения -- мы ещё ничего не записали в shadow map
\end{itemize}
\end{frame}

\begin{frame}[fragile]
\frametitle{Задание 3}
\slideimage{3.png}
\end{frame}

\begin{frame}[fragile]
\frametitle{Задание 4}
Генерируем shadow map
\begin{itemize}
\item Создаём новый шейдер для генерации теней: он просто записывает \mintinline{wgsl}|position| входной вершины как \mintinline{wgsl}|@builtin(position)| выходного значения шейдера (другие атрибуты ему не нужны)
\item Фрагментный шейдер тут тоже не нужен :)
\item Создаём новый render pipeline для генерации теней: он использует новый шейдер, те же настройки depth-stencil, но не выводит никакой цвет (fragment state = null)
\item Bind группы ему тоже не нужны
\item Заводим новый render pass \textbf{до} основного: у него есть текстура shadow map в качестве depth attachment (\mintinline{cpp}|store_op = Store|), но нет ни одного color attachment
\item В этом render pass рисуем нашу 3D модель используя новый render pipeline (\mintinline{cpp}|set_vertex_buffer()|, \mintinline{cpp}|set_index_buffer()|, \mintinline{cpp}|draw_indexed()|)
\item Модель должна появиться в дебажном квадрате
\end{itemize}
\end{frame}

\begin{frame}[fragile]
\frametitle{Задание 4}
\slideimage{4.png}
\end{frame}

\begin{frame}[fragile]
\frametitle{Задание 5}
Вычисляем проекцию для shadow map
\begin{itemize}
\item Заводим оси проекции:
\begin{itemize}
\item \mintinline{cpp}|Z = -light_direction|
\item \mintinline{cpp}|X| и \mintinline{cpp}|Y| -- любые, образующие с \mintinline{cpp}|Z| ортонормированный базис
\end{itemize}
\item Эти векторы дадут \textbf{строки} матрицы перевода в систему координат теней
\item Диапазон значений Z -- \([0, 1]\), поэтому тени обрежутся плоскостью, проходящей через начало координат
\item Это проще всего исправить, домножив проекцию теней на матрицу преобразования \(z \mapsto \frac{z+1}{2}\)
\item Получившуюся матрицу нужно передать в шейдер генерации теней (проще всего как immediate) и применить к входным вершинам
\item Shadow map в дебажном прямоугольнике должна начать двигаться со временем
\end{itemize}
\end{frame}

\begin{frame}[fragile]
\frametitle{Задание 5}
\slideimage{5.png}
\end{frame}

\begin{frame}[fragile]
\frametitle{Задание 6}
Накладываем тени
\begin{itemize}
\item В lights uniforms добавляем матрицу shadow projection, вычисленную на предыдущем шаге
\item В основной шейдер добавляем bindings для shadow map и sampler
\item Во фрагментном шейдере вычисляем, находится ли в тени данный пиксель (см. слайды лекции)
\item Применяем освещение от источника света только если мы \textbf{не находимся} в тени (ambient применяем всегда!)
\item \textbf{NB:} Будет очень сильный shadow acne
\end{itemize}
\end{frame}

\begin{frame}[fragile]
\frametitle{Задание 6}
\slideimage{6.png}
\end{frame}

\begin{frame}[fragile]
\frametitle{Задание 7}
Добавляем shadow bias
\begin{itemize}
\item Вычитаем маленькое значение из \mintinline{cpp}|shadow_coord.z|
\item Значение нужно подобрать, чтобы исчез shadow acne, но чтобы тени не съехали слишком сильно
\end{itemize}
\end{frame}

\begin{frame}[fragile]
\frametitle{Задание 7}
\slideimage{7.png}
\end{frame}

\begin{frame}[fragile]
\frametitle{Задание 8*}
Добавляем PCF
\begin{itemize}
\item Создаём ещё один sampler, \mintinline{cpp}|filter = Linear|, \mintinline{cpp}|compare = Less| (старый sampler нельзя переиспользовать -- сломается дебажная визуализация shadow map)
\item Добавляем его в lights bind group, \mintinline{cpp}|sampler = Comparison|
\item Добавляем соответствующий binding в основной шейдере, тип sampler'а -- \mintinline{wgsl}|sampler_comparison|
\item При наложении теней в основном шейдере используем не \mintinline{wgsl}|textureSample|, а \mintinline{wgsl}|textureSampleCompare| -- она вернёт \textbf{множитель}, на который нужно домножить освещение от источника света
\item Вблизи тени должны стать гладкими
\end{itemize}
\end{frame}

\begin{frame}[fragile]
\frametitle{Задание 8*}
\slideimage{8.png}
\end{frame}

\end{document}
